% F -- The whole report on one page: three findings, three numbers.
% NB: 'cap' is a reserved TikZ key (line cap); the caption style is 'capt'.
\begin{figure}[!htbp]
\centering
\fitwidth{%
\begin{tikzpicture}[font=\small,
  card/.style={anchor=north west,draw=clrule,fill=white,rounded corners=3pt,
               inner sep=0pt,minimum width=4.70cm,minimum height=6.65cm},
  eyeb/.style={anchor=north west,align=left,text width=4.20cm,
               font=\tiny\bfseries,text=clrule},
  big/.style={anchor=north west,font=\fontsize{25}{25}\selectfont\bfseries},
  capt/.style={anchor=north west,align=left,text width=4.20cm,font=\scriptsize,
               text=clmodel},
  mic/.style={font=\tiny,text=clrule},
]
% ============================ card 1 ====================================
\begin{scope}
\node[card] at (0,0) {};
\node[eyeb] at (0.25,-0.22) {THE MODEL BEATS THE ONE-LINE RULE BY};
\node[big,text=clgain] at (0.25,-0.92) {+4.3};
\node[capt] at (0.25,-1.98)
  {more correct answers per hundred than \emph{``the same antenna as a moment
   ago''}: 54.5\,\% against 50.2\,\% on a thousand people never seen during
   training.};
\fill[clbase] (0.25,-4.34) rectangle (1.56,-4.10);
\fill[clgain] (1.56,-4.34) rectangle (1.67,-4.10);
\node[mic,anchor=west,text=clgain] at (1.77,-4.22) {54.5};
\node[mic,anchor=north west] at (0.25,-4.60)
  {blue: the rule; orange: the model's gain};
\node[capt] at (0.25,-5.35)
  {What produced it was not a bigger model, but sorting people by \emph{when}
   they move.};
\end{scope}

% ============================ card 2 ====================================
\begin{scope}[shift={(5.05,0)}]
\node[card] at (0,0) {};
\node[eyeb] at (0.25,-0.22) {MORE PEOPLE STOP PAYING AT};
\node[big,text=clmodel] at (0.25,-0.92) {1\,900};
\node[capt] at (0.25,-1.98)
  {training people. Below that the gain climbs from 0.6 to 4.2. Above it, a
   ten-fold increase to 21{,}000 buys \textbf{0.1}, and zero is inside every
   interval.};
\draw[clrule] (0.25,-4.85) -- (4.15,-4.85);
\draw[clmodel,very thick]
  (0.25,-4.74) -- (1.03,-4.35) -- (1.45,-4.22) -- (1.75,-4.09)
  -- (2.67,-4.08) -- (3.34,-4.06) -- (3.73,-4.00) -- (4.05,-4.08);
\draw[clloss,dashed] (1.75,-4.85) -- (1.75,-3.90);
\node[mic,anchor=south,text=clloss] at (1.75,-3.88) {1{,}900};
\node[mic,anchor=north west] at (0.25,-4.94)
  {what the model gains, against training people};
\node[capt] at (0.25,-5.35)
  {Two thousand people is the operating point; collecting ten times more is
   not where the next gain lies.};
\end{scope}

% ============================ card 3 ====================================
\begin{scope}[shift={(10.10,0)}]
\node[card] at (0,0) {};
\node[eyeb] at (0.25,-0.22) {THE CALMEST AND THE MOST RESTLESS KIND OF USER ARE APART BY};
\node[big,text=clGthree] at (0.25,-0.92) {22.4};
\node[capt] at (0.25,-1.98)
  {correct answers per hundred. Over the whole range of training sizes, from 400
   people to 21{,}000, the gain moves by 0.4.};
\fill[clGthree] (0.25,-4.29) rectangle (1.95,-4.05);
\node[mic,anchor=west,text=clGthree] at (2.05,-4.17) {22.4 between kinds of user};
\fill[clnonsig] (0.25,-4.79) rectangle (0.29,-4.55);
\node[mic,anchor=west] at (0.39,-4.67) {0.4 between training sizes};
\node[capt] at (0.25,-5.35)
  {Who the person is decides the score; how much data was collected barely
   moves it.};
\end{scope}

% ============================ footer ====================================
\node[anchor=north west,draw=clbase,dashed,fill=clbase!8,rounded corners=3pt,
      inner sep=7pt,align=left,text width=14.35cm,font=\footnotesize]
      at (0,-6.95)
  {\textbf{And the three largest single improvements were about measuring, not
   modelling.} Reading the best saved copy of a run instead of the last one was
   worth \textbf{15.5} correct answers per hundred. Recognising that where the
   day is cut belongs to the ruler retired a \textbf{1.1} ``result'' that was
   noise. And choosing a set of people on which the trivial rule does not already
   score 69\,\% is what made the question measurable at all.};
\end{tikzpicture}}
\caption{The whole report in three numbers}
\label{fig:results-card}
\figtext{%
\figpara{Summary of key findings in Figure~\ref{fig:results-card} Cards 1--3.}{Each card in Figure~\ref{fig:results-card} is the answer to one of the three
questions the internship actually asked: does a retrained language model beat the
trivial rule (Card 1), how much data does that take (Card 2), and where does the remaining
difficulty live (Card 3). They are ordered by how firmly they are established. The
\textbf{+4.3} in Card 1 is one measurement on one set of people, repeated once on a second
group; the \textbf{1{,}900} in Card 2 rests on eight training sizes and a resampling test;
the \textbf{22.4} in Card 3 is the most robust of the three, because it appears in every
chapter that split its results by person.}

\figpara{Methodological corrections highlighted in Figure~\ref{fig:results-card}.}{In Figure~\ref{fig:results-card}, on a task whose honest yardstick is
one line of code, three corrections to how things were measured were each worth
more than any modelling idea tested. That is where a new intern taking over this project should look first.}}

\figsource{\nbfinal}{train\_cellid\_final.ipynb}{with every number traced to its run in Table~\ref{tab:provenance}}
\end{figure}
